% !TeX root = ../动手学强化学习笔记.tex
% 由主文档通过 \input 载入；请编译主文档。
\setcounter{chapter}{9} % 保留原书章号。
\chapter{离策略学习与重要性采样}
\label{ch:rl-offpolicy}

离策略学习允许使用另一种行为产生的数据，但“数据可以复用”不意味着“任何数据都足够”。首先必须区分生成数据的行为策略（behavior policy）$b$ 与待评价的目标策略（target policy）$\pi$，再检查目标策略需要的行为是否被数据覆盖。

\section{覆盖条件与分布校正}

对给定状态 $s$，若 $\pi(a\mid s)>0$ 的动作均满足 $b(a\mid s)>0$，则任意可积函数 $f$ 满足
\[
\sum_a\pi(a\mid s)f(a)=\sum_a b(a\mid s)\frac{\pi(a\mid s)}{b(a\mid s)}f(a).
\]
这就是重要性采样（importance sampling）：来自低概率行为动作的样本需要较大权重。若行为策略从不采取某目标动作，就没有相应样本能够被重新加权。

在共同初始分布和环境动力学下，固定策略的轨迹概率都含相同的初始分布与转移核，取比后这些项抵消。定义
\[
\rho_t=\frac{\pi(A_t\mid S_t)}{b(A_t\mid S_t)},\qquad
\rho_{t:u}=\prod_{k=t}^{u}\rho_k.
\]
空乘积定义为一。在目标轨迹分布相对于行为轨迹分布绝对连续、回报可积的条件下，给定起始状态的完整回报满足 $v_\pi(s)=\mathbb E_b[\rho_{0:T-1}G_0\mid S_0=s]$。若初始分布不同，还需要初始密度比；估计 $q_\pi(s,a)$ 且第一步动作已固定时，策略比从下一动作开始计算。

\section{普通与加权重要性采样}

设 $n$ 个独立回合给出回报 $G_i$ 与完整权重 $W_i$。普通估计和加权估计分别为
\[
\widehat v_{\rm OIS}=\frac1n\sum_{i=1}^nW_iG_i,\qquad
\widehat v_{\rm WIS}=\frac{\sum_{i=1}^nW_iG_i}{\sum_{i=1}^nW_i}.
\]
普通估计在上述条件下无偏，有限方差还要求 $\mathbb E_b[W^2G^2]<\infty$；加权估计是比值，一般有限样本有偏，但当 $\mathbb E_bW=1$、$\mathbb E_b|WG|<\infty$ 时由强大数定律得到一致性。分母为零时加权估计尚未定义，不能把零当成有意义的价值估计。

\begin{example}[概率比的乘积]
目标策略和行为策略对轨迹中两次已选动作的概率分别为 $(0.6,0.2)$ 与 $(0.3,0.4)$。两步概率比为 $2$ 和 $1/2$，完整权重为 $1$；这并不意味着中间每项奖励都应使用权重一。若两个独立回合的 $(W,G)$ 为 $(2,3)$ 与 $(1/2,1)$，普通估计为 $3.25$，加权估计为 $2.6$。
\end{example}

长轨迹中很多概率比相乘可能产生极端权重。可记录描述性有效样本量 $(\sum_iW_i)^2/\sum_iW_i^2$；它能提示权重集中，但不是任意情况下精确的独立样本数或置信保证。

\section{逐决策校正与多步学习}

奖励 $R_{k+1}$ 只受此前动作影响，不依赖尚未执行的未来动作。因此在有限回合、相应可积条件下，逐决策重要性采样给出
\[
v_\pi(s)=\mathbb E_b\!\left[\sum_{k=0}^{T-1}\gamma^k\rho_{0:k}R_{k+1}\,\middle|\,S_0=s\right].
\]
固定有限时域下，可对每一项奖励的轨迹前缀单独作换测度证明，再交换有限求和与期望。随机回合长度则需可积性保证相应交换。

多步自举同时包含真实奖励与后继价值，校正时应分别核对每项依赖哪些动作。简单截断权重 $\min(c,W)$ 往往降低极端样本影响，但改变了期望目标；除非配合专门修正，不能仍称其无偏。树备份（tree-backup）用目标策略概率对未采样分支的价值估计取期望，以更多自举代替部分高方差比值；这种取舍同样依赖价值估计的准确性。

\section{为什么 Q-learning 不乘完整轨迹比}

Q-learning 更新给定 $(s,a)$ 后的条件期望，其一步目标为 $R+\gamma\max_{a'}Q(S',a')$。行为策略决定哪些 $(s,a)$ 被访问，但在相同环境中，给定这一对后的 $(R,S')$ 条件分布仍是正确转移核。因此表格算法不必对完整轨迹加权，只需充分访问各对并满足相应收敛条件。

函数近似时，不同状态动作对共享参数，行为分布改变也会改变总体更新权重；上述条件分布正确不再足以保证参数更新稳定。离线学习还无法主动补充缺失动作，因此“Q-learning 是离策略的”不等于“直接在任意固定数据上训练就可靠”。

\section{习题与实验}
\begin{enumerate}
\item 两动作行为概率为 $(0.8,0.2)$，目标概率为 $(0.4,0.6)$，动作收益为 $(1,3)$。分别直接计算目标收益和重要性加权的行为期望。
\item 构造一个动作从未被行为策略采取的例子，说明即使获得无限多条行为轨迹，仍可能无法确定目标价值。
\item 在固定目标策略下改变行为策略与回合长度，比较普通、加权和逐决策估计的误差及权重分布。不要用相同训练数据反复筛选目标策略后仍宣称原有无偏评价结论。
\end{enumerate}





